\chapter{Introducción: Búsquedas en IA}

\section{Formulación}
La Inteligencia Artificial Clásica trata entornos con las siguientes características:
\begin{itemize}
    \item Estático: El entorno no varía.
    \item Observable: Conocemos todas las variables que nos interesan.
    \item Discreto: Los valores tienen límite de división.
    \item Determinístico: El estado siguiente es consecuencia del estado actual y de la acción realizada.
\end{itemize}

\section{Características}
Se puede construir un grafo con la representación del problema. Las características de este grafo serán:
\begin{itemize}
    \item Cada estado se representará mediante un nodo.
    \item Deberá tener un estado inicial y, al menos, un estado final.
    \item Las aristas representan las acciones entre los estados.
    \item Puede ser cíclico.
\end{itemize}

\section{Espacio De Estados}
\subsection{Árbol De Búsqueda}

\section{Solución}

\section{Ejemplos}

\section{Búsqueda De Soluciones}

\section{Algoritmo General De Búsqueda}

\section{Estrategias De Búsqueda}

\section{Criterios Para Evaluar Las Estrategias}

\section{Búsquedas Con Adversarios}



\section{Búsqueda En El Espacio De Estados (Temp)}


\subsection{Estructuras de Datos y Definiciones}
A continuación se define la estructura necesaria para un algoritmo de búsqueda genérico (como A* o Dijkstra).

\begin{lstlisting}[caption=Interfaces y Clase Nodo]
public interface Estado {
    public Estado aplicarAccion(Action a);
    public boolean esMeta();
    public boolean esValido();
    public int h(); // Heuristica: coste estimado hasta la meta
}

public interface Action {
    public int costeAccion();
}

public class Nodo implements Comparable<Nodo> {
    private int g;            // Coste acumulado (inicio -> actual)
    private Estado estado;
    private Nodo padre;
    private Action accion;

    public Nodo(int g, Estado estado, Nodo padre, Action accion) {
        this.g = g;
        this.estado = estado;
        this.padre = padre;
        this.accion = accion;
    }

    public int getF() { return this.g + this.estado.h(); }
    public int getG() { return g; }
    public Estado getEstado() { return estado; }
    public Nodo getPadre() { return padre; }
    public Action getAccion() { return accion; }

    @Override
    public int compareTo(Nodo otro) {
        return Integer.compare(this.getF(), otro.getF());
    }
}
\end{lstlisting}

\subsection{Algoritmo de Búsqueda General}
El siguiente pseudocódigo implementa una búsqueda que evita ciclos mediante un conjunto de estados cerrados y utiliza una cola de prioridad para optimizar la expansión.

\begin{lstlisting}[caption=Algoritmo de Búsqueda corregido]
public class BusquedaEspacioEstados {
    public List<Action> buscar(Estado inicial, List<Action> acciones) {
        PriorityQueue<Nodo> abiertos = new PriorityQueue<>();
        Set<Estado> cerrados = new HashSet<>();

        abiertos.add(new Nodo(0, inicial, null, null));

        while (!abiertos.isEmpty()) {
            Nodo actual = abiertos.poll();
            Estado estadoActual = actual.getEstado();

            if (estadoActual.esMeta()) {
                return reconstruirCamino(actual);
            }

            if (!cerrados.contains(estadoActual)) {
                cerrados.add(estadoActual);

                for (Action a : acciones) {
                    Estado sigEstado = estadoActual.aplicarAccion(a);
                    if (sigEstado != null && sigEstado.esValido()) {
                        int nuevoG = actual.getG() + a.costeAccion();
                        abiertos.add(new Nodo(nuevoG, sigEstado, actual, a));
                    }
                }
            }
        }
        return null; // No hay solucion
    }

    private List<Action> reconstruirCamino(Nodo nodo) {
        List<Action> camino = new ArrayList<>();
        Nodo aux = nodo;
        while (aux != null && aux.getAccion() != null) {
            camino.add(0, aux.getAccion());
            aux = aux.getPadre();
        }
        return camino;
    }
}
\end{lstlisting}

\subsection{Notas de Implementación}
\begin{itemize}
    \item \textbf{Variantes del algoritmo}:
    \begin{itemize}
        \item \textbf{A*}: Se activa cuando \texttt{h()} devuelve una estimación optimista del coste restante.
        \item \textbf{Dijkstra}: Se obtiene si la función heurística \texttt{h()} devuelve siempre 0.
        \item \textbf{BFS}: Se obtiene si el coste de toda acción es 1 y \texttt{h()} es 0.
    \end{itemize}
    \item \textbf{Control de Ciclos}: El uso de \texttt{Set<Estado> cerrados} es vital para evitar bucles infinitos en grafos.
    \item \textbf{Requisitos de Java}: Para que la detección de estados repetidos funcione correctamente, la clase que implemente \texttt{Estado} debe sobreescribir obligatoriamente los métodos \texttt{equals()} y \texttt{hashCode()}.
\end{itemize}

% ---------Bibliografía del Tema------------
\bigskip
\noindent\rule{\textwidth}{0.4pt} 

\section*{Bibliografía del Tema}
\addcontentsline{toc}{section}{Bibliografía del Tema}

\begin{itemize}
    \item \textbf{Russell, S. \& Norvig, P.} (2003). \textit{Artificial Intelligence: A Modern Approach} (2nd ed.). Prentice Hall. Chapter 3.
    
\end{itemize}
\newpage